\documentclass[11pt]{article}
\usepackage[margin=1.1in]{geometry}
\usepackage{amsmath,amssymb}
\usepackage{booktabs}
\usepackage{array}
\usepackage[hidelinks]{hyperref}
\usepackage{xcolor}
\usepackage{parskip}
\usepackage{graphicx}
\usepackage{tikz}

% ---- reversible-gate drawing primitives (control pins + active pin) ----
\newcommand{\pc}[2]{\fill (#1,#2) circle (2.6pt);}                       % positive control
\newcommand{\nc}[2]{\draw[fill=white] (#1,#2) circle (2.6pt);}          % negative control
\newcommand{\tg}[2]{\draw (#1,#2) circle (4.2pt);                        % target (XOR)
  \draw (#1,#2) ++(-4.2pt,0) -- ++(8.4pt,0);
  \draw (#1,#2) ++(0,-4.2pt) -- ++(0,8.4pt);}
\newcommand{\tgb}[2]{\tg{#1}{#2}                                         % complemented target (g57)
  \draw[line width=1pt] (#1,#2) ++(-4.8pt,8pt) -- ++(9.6pt,0);}
\newcommand{\gcol}[3]{\draw (#1,#2)--(#1,#3);}                           % gate's vertical line
\newcommand{\wires}[1]{% #1 = number of gate columns
  \foreach \yy/\lab in {0/{$a^\ast$},-1/{$b_0$},-2/{$b_1$},-3/{$c_0$},-4/{$c_1$}}{%
    \draw (0.55,\yy)--({#1+0.45},\yy); \node[left] at (0.5,\yy){\small\lab};}
  \foreach \j in {1,...,#1}{\node at (\j,0.82){\scriptsize\j};}}
\newenvironment{cgcirc}{\begin{tikzpicture}[x=0.85cm,y=0.8cm,line width=0.7pt]}
  {\end{tikzpicture}}

\newcommand{\CNOT}{\textsc{cnot}}
\newcommand{\CCNOT}{\textsc{ccnot}}
\newcommand{\code}[1]{\texttt{#1}}
\newcommand{\gf}{\mathrm{GF}(2)}
\newcommand{\pe}{\mathrel{{+}{=}}}    % t += e  means  t <- t XOR e over GF(2)
\newcommand{\at}{a^{\ast}}            % the targeted carrier of A

\title{\textbf{Nonlinear Re-Randomization and CG Diversification}\\[2pt]
\large Why the new gadgetize reverted to the nonlinear RGs, and why each CG
is now drawn from a random menu}
\author{local\_mixing project}
\date{July 20, 2026 \\ \small branch \code{ssg-gen-mix-clean}, commits
\code{813bf256}, \code{4fe82458}}

\begin{document}
\maketitle

\begin{abstract}
\noindent The 2026-07-20b revision of the new-gadgetize recipe makes two
changes to the gadget body: the linear \CNOT{} re-randomization gadgets
$\{RG2, RG3\}$ are replaced by the legacy \emph{nonlinear} g57 networks
$\{RG1, RG2, RG3\}$ (one uniform draw per shared gate), and the fixed
four-fragment realization of the shared computation gadget (CG) is replaced
by a per-gate uniform draw over \emph{seven} structurally different
realizations. This note records the motivation: what the degree-bounded
reconstruction experiments showed about the linear design, what the
reversal buys and what it deliberately gives up, and why implementation
diversity of the CG is itself a security measure. Companion changes (the
final commuting-order rerandomization and the sandwich column float) are
summarized at the end.
\end{abstract}

\section{Setting}

The two-share gadget maps a circuit $C$ on $n$ wires to an equivalent
circuit on $2n$ wires. Every logical value $v$ is carried as an XOR pair
$v = s_v \oplus p_v$; the construction consists of random Z bookends on the
aux half, a $W_i$ encode/decode curtain of \CNOT{}s that installs and
removes the sharing, one \emph{computation gadget} (CG) per source gate,
and a stream of \emph{re-randomization gadgets} (RGs) that churn the
sharing between CGs. The zero-slice \CNOT/\CCNOT{} preblock guards the
input port.

The original new-gadgetize design chose \emph{linear} RGs (a 3-\CNOT{}
re-pairing RG2 and a 2-\CNOT{} mask-refresh RG3) and one fixed
four-fragment CG. Both choices optimized a specific adversary model:
\emph{gate-local non-completeness} --- no single physical gate ever reads
both carriers of one value, so every prefix of the circuit is first-order
probe-masked --- plus gate-count leanness.

\section{The adversary that changed the decision}

The project's security target is indistinguishability against \emph{all
feasible distinguishers}, evaluated against concrete cheap attack
families. The family that drove this revision is \emph{degree-bounded
algebraic reconstruction} (\code{hmap\_affine}): for prefixes $C_i$ of the
original and $G_j$ of the gadgetized circuit, fit each bit of $C_i(x)$ as
a degree-$d$ $\gf$ function of \emph{all} of $G_j$'s wires (train/holdout
split; inconsistency scores $0.5$). Low error means the original
computation's progress is legible at that point of the gadget; the ideal
map is flat at $0.5$ except the two forced I/O corners. The measure is
invariant to affine re-encodings, which is exactly what makes it the right
lens for a sharing-based gadget: \emph{the XOR sharing itself is free for
this adversary}.

Two facts came out of the measurement campaign:

\begin{enumerate}
\item \textbf{Degree-2 reconstruction is more persistent than degree-1
under mixing.} Mixing erodes the degree-1 (affine) diagonal quickly, but
the degree-2 diagonal survives much longer. The affine conjugation twists
of the mixing stage cannot address this: an affine change of variables is
transparent at every degree.
\item \textbf{The gadget encodings differ sharply.} The old (legacy g57)
gadget's output is \emph{not} recoverable even by restricted degree-2
predictors, while the linear-RG \CNOT{} gadget is affine-transparent at
its interface. Computing the ANF degree of each component isolated the
cause: the sharing setup ($M$, $W_i$) is affine everywhere; the
nonlinearity of the legacy gadget lives in its gate application (degree
2, once per source gate, compounding) and its re-randomizers --- RG1
(degree 3), RG2 (degree 2), RG3 (degree 2). The linear replacement kept
the setup and \emph{removed exactly the nonlinear part}.
\end{enumerate}

The structural point is simple: affine maps compose to affine maps. With
linear RGs, the entire re-randomization stream between any two CGs is one
affine transformation of the wire state, so a degree-$d$ predictor sees
through an \emph{arbitrarily long} RG stream at no cost in degree. The
stream churned the sharing but hid nothing from the adversary that
matters. Nonlinear RGs break this composition: every RG multiplies
algebraic degree into the transition, and reconstruction across a stretch
of the body must now pay degree for every RG it crosses.

\section{The reversal: nonlinear $\{RG1, RG2, RG3\}$, one per CG}

The revision reinstates the legacy g57 networks, drawn uniformly at one RG
per CG gap (\code{--rg-frequency}, default 1 on this path):

\begin{center}
\begin{tabular}{lllc}
\toprule
RG & effect on the sharing & gates & ANF degree \\
\midrule
RG1 & swap the virtual values of two pairs & 6 g57 & 3 \\
RG2 & re-pair two pairs crosswise & 6 g57 & 2 \\
RG3 & XOR $r_1 \vee \lnot r_2$ of two foreign wires into both carriers & 2 g57 & 2 \\
\bottomrule
\end{tabular}
\end{center}

\textbf{What is bought.} Nonlinear transitions of degree 2--3 that compound
along the body; a re-randomization stream that low-degree reconstruction
cannot cross for free; and, as a side effect, nonlinear material (the RG3
reads of two \emph{random} foreign wires) distributed through the entire
body rather than confined to the bookends.

\textbf{What is paid, deliberately.} RG1 and RG2 gates read both carriers
of a value, so gate-local non-completeness --- first-order probe masking
--- is given up. This is a conscious ranking of adversaries: the probing
model was a secondary concern, while degree-bounded reconstruction is a
demonstrated, cheap, working attack on the linear design. (RG3 keeps
single-carrier reads; the masking-safe CG variants below keep the option
of a partially masked body alive.)

\textbf{Measured effect} ($n{=}16$, pre-mixing, old vs.\ revised recipe):
the maps flatten (degree-1 $\sigma$ $0.125 \to 0.106$, degree-2 $0.161 \to
0.149$; the interior column-mean dip that leaked the temporal layout
largely evens out), and the best-case interior degree-1 reconstruction
error roughly doubles ($\approx 0.05$--$0.15 \to 0.2$--$0.3$). One honest
caveat: the degree-2 \emph{snapshot} diagonal persists in both versions.
This is inherent to the linear XOR sharing --- at any prefix the logical
value \emph{is} $s \oplus p$, a degree-1 function of the wires ---
independent of how nonlinear the transitions are. Erasing snapshots is the
mixing stage's job (structural DB/split mixing is what erodes it), or, in
a future design, a nonlinear share \emph{encoding}.

\section{CG diversification: a seven-variant menu}

The CG applies $A \mathrel{\oplus}= B \vee \lnot C$ over the shares, i.e.\
$f = 1 \oplus C \oplus BC$ with $B = b_0 \oplus b_1$, $C = c_0 \oplus
c_1$. The original design emitted one fixed shape --- four mixed-polarity
conjunction fragments --- for every source gate.

\textbf{Why fixed is bad.} A fixed CG is a syntactic beacon: a fixed
period, a fixed local gate-type pattern, and a fragment census that counts
the source gates exactly. Regularity is precisely what static profiling
and SAT-style attacks feed on --- a solver that learns one CG shape has
learned them all. And the fragments themselves are cold material for the
mixing stage: the frozen replacement store matches g57 windows at
${\sim}99\%$ but mpmct fragment windows at ${\sim}2.8\%$, so a
fragment-heavy body resists exactly the phase-A mixing that is supposed
to dissolve it.

\textbf{The menu.} Each source g57 is now emitted as one of seven
realizations, drawn uniformly per gate, with the carrier roles (which
carrier of $A$ is targeted; share/pad order of $B$ and of $C$) also
randomized per draw:

\begin{center}
\small
\begin{tabular}{cllp{5.6cm}}
\toprule
\# & name & shape & idea \\
\midrule
0 & collapse both & 4 \CNOT{} + 1 g57 & $b_0{\oplus}{=}b_1$, $c_0{\oplus}{=}c_1$ put $B$, $C$ on single wires; one g57 fires $f$; restore \\
1 & collapse $c$ & 2 \CNOT{} + 2 g57 + $\lnot$\CNOT{} & pair sums to $BC$; $a \mathrel{\oplus}= \lnot c_0$ adds $1 \oplus C$ \\
2 & collapse $b$ & 2 \CNOT{} + 1 g57 + 1 frag & $[a,B,c_0]$ plus fragment $\lnot B \wedge c_1$ \\
3 & linear tail & $\lnot$\CNOT{} + \CNOT{} + 4 g57 & no collapse; quad $[a,b_i,c_j]$ sums to $BC$; masking-safe \\
4 & legacy GADGET & 6 g57 & the classic nonlinear network \\
5 & quad + pair & 5 g57 + 1 frag & quad $BC$, then $(\lnot c_0 \wedge c_1) + [a,c_1,c_0] = 1 \oplus C$ \\
6 & 4-cube ESOP & 4 fragments & the previous fixed CG; masking-safe \\
\bottomrule
\end{tabular}
\end{center}

Vocabulary is restricted to g57s and pure conjunctions with one or two
controls of any polarity. In particular, \emph{no bare X gates}: several
variants algebraically need a ``$\oplus 1$'', but an X census would count
the source gates, so every constant is folded into a negative-control
\CNOT{} or a complemented fragment.

\textbf{Gate-by-gate.} Each realization writes only the target carrier
$\at$ (one of $A$'s two carriers, chosen per draw); $b_0, b_1$ and $c_0,
c_1$ are the two carriers of $B$ and of $C$ in a per-draw random order.
Notation, all over $\gf$ (so $+$ is $\oplus$ and $\overline{w} = 1 \oplus
w$): $t \pe w$ is a \CNOT{} ($t \leftarrow t \oplus w$); $t \pe
\overline{w}$ a negative-control \CNOT{}; $t \pe (x \vee \overline{y})$ a
g57; and $t \pe (\ell\,\ell')$ a two-control conjunction fragment
(juxtaposition = AND, each literal $w$ or $\overline{w}$). Recall $f = 1
\oplus C \oplus BC = B \vee \overline{C}$. Gates execute top to bottom. The
same seven circuits are drawn in standard reversible-gate notation in
Appendix~\ref{app:cgfigs}.

\newcommand{\note}[1]{\hfill{\footnotesize\itshape #1}}
\newcommand{\cgvar}[1]{\smallskip\noindent\textbf{#1}\par\nopagebreak}

\cgvar{0 --- collapse both} (4 \CNOT{} + 1 g57):
\begin{enumerate}\itemsep1pt
\item $b_0 \pe b_1$ \note{now $b_0 = B$}
\item $c_0 \pe c_1$ \note{now $c_0 = C$}
\item $\at \pe (b_0 \vee \overline{c_0})$ \note{$= B \vee \overline{C} = f$}
\item $c_0 \pe c_1$ \note{restore $c_0$}
\item $b_0 \pe b_1$ \note{restore $b_0$}
\end{enumerate}

\cgvar{1 --- collapse $c$} (2 \CNOT{}-type + 2 g57):
\begin{enumerate}\itemsep1pt
\item $c_0 \pe c_1$ \note{now $c_0 = C$}
\item $\at \pe (b_0 \vee \overline{c_0})$
\item $\at \pe (b_1 \vee \overline{c_0})$ \note{the two g57s sum to $BC$}
\item $\at \pe \overline{c_0}$ \note{adds $\overline{C} = 1 \oplus C$}
\item $c_0 \pe c_1$ \note{restore $c_0$}
\end{enumerate}

\cgvar{2 --- collapse $b$} (2 \CNOT{} + 1 g57 + 1 fragment):
\begin{enumerate}\itemsep1pt
\item $b_0 \pe b_1$ \note{now $b_0 = B$}
\item $\at \pe (b_0 \vee \overline{c_0})$ \note{$= 1 \oplus c_0 \oplus B c_0$}
\item $\at \pe (\overline{b_0}\, c_1)$ \note{$= c_1 \oplus B c_1$; total $= f$}
\item $b_0 \pe b_1$ \note{restore $b_0$}
\end{enumerate}

\cgvar{3 --- linear tail} ($\overline{\CNOT{}}$ + \CNOT{} + 4 g57; masking-safe, no carrier collapsed):
\begin{enumerate}\itemsep1pt
\item $\at \pe \overline{c_0}$
\item $\at \pe c_1$ \note{gates 1--2 give $\overline{c_0} \oplus c_1 = 1 \oplus C$}
\item $\at \pe (b_0 \vee \overline{c_0})$
\item $\at \pe (b_1 \vee \overline{c_0})$
\item $\at \pe (b_0 \vee \overline{c_1})$
\item $\at \pe (b_1 \vee \overline{c_1})$ \note{the four g57s sum to $BC$}
\end{enumerate}

\cgvar{4 --- legacy GADGET} (6 g57; $c_1$ borrowed as restored scratch):
\begin{enumerate}\itemsep1pt
\item $c_1 \pe (b_0 \vee \overline{b_1})$ \note{scratch on $c_1$}
\item $\at \pe (c_1 \vee \overline{b_1})$
\item $\at \pe (b_0 \vee \overline{c_1})$
\item $c_1 \pe (b_0 \vee \overline{b_1})$ \note{restore $c_1$}
\item $\at \pe (b_1 \vee \overline{c_0})$
\item $\at \pe (c_0 \vee \overline{b_0})$
\end{enumerate}

\cgvar{5 --- quad $+$ pair} (5 g57 + 1 fragment; no carrier collapsed):
\begin{enumerate}\itemsep1pt
\item $\at \pe (b_0 \vee \overline{c_0})$
\item $\at \pe (b_1 \vee \overline{c_0})$
\item $\at \pe (b_0 \vee \overline{c_1})$
\item $\at \pe (b_1 \vee \overline{c_1})$ \note{the four g57s sum to $BC$}
\item $\at \pe (\overline{c_0}\, c_1)$
\item $\at \pe (c_1 \vee \overline{c_0})$ \note{gates 5--6 sum to $1 \oplus C$}
\end{enumerate}

\cgvar{6 --- 4-cube ESOP} (4 fragments; the previous fixed CG, masking-safe):
\begin{enumerate}\itemsep1pt
\item $\at \pe (\overline{c_1}\, b_0)$
\item $\at \pe (c_1\, \overline{b_1})$
\item $\at \pe (c_0\, b_1)$
\item $\at \pe (\overline{c_0}\, \overline{b_0})$ \note{the four cubes XOR to $B \vee \overline{C}$}
\end{enumerate}

\smallskip
Every non-target carrier that a variant touches (the collapses in 0--2, the
scratch in 4) is written an even number of times and so is restored exactly
at the wire level; verified exhaustively over all $7$ variants, $8$ carrier
role assignments, and $64$ inputs.

\textbf{What diversification buys.}
\begin{itemize}
\item \emph{Irregularity}: no fixed CG period, no repeated shape, varying
local gate-type statistics --- less structure for syntactic profiling and
for solvers to latch onto. The carrier-role randomization scrambles pin
patterns even between same-variant draws.
\item \emph{DB-warmth}: the body census flips from fragment-heavy to
g57-dominant (at $n{=}16$, g57 / \CNOT{} / fragment counts $822/277/707
\to 1269/486/264$), with \CNOT{}s only sprinkled --- and a few \CNOT{}s
within many g57s are well tolerated by the replacement store, while the
g57 mass mixes at the ${\sim}99\%$ match rate.
\item \emph{More nonlinearity in the body}: five of seven variants carry
g57 material into the CG itself, compounding with the nonlinear RG stream.
\end{itemize}

Cost: the menu averages ${\sim}5.1$ gates per CG versus 4, i.e.\
${\sim}12\%$ body growth at $n{=}16$ --- accepted.

Every variant is verified exhaustively (all $7$ variants $\times$ $8$ role
draws $\times$ $64$ carrier assignments: exact update, wire-level
restoration of every non-target carrier, vocabulary compliance), and the
driver's pinned zero-slice functionality check runs on every build.

\section{Companion changes, briefly}

\textbf{Commuting-order rerandomization} (\code{commuting\_shuffle},
\code{5689c8ea}). The construction-time layout $S_1\,|\,Z\,|\,W\,|\,
\text{body}\,|\,W^{-1}\,|\,Z$ was legible in the heatmaps as an S-shape
--- no mixing between the computation block and the bookends, with the
affine $W_i$ curtain sitting between them. The emitted order is now a
fresh random order that preserves only the relative order of gate pairs
that \emph{provably} collide (\code{XGate::collides}: ``commute unless
proven otherwise'', including the opposite-polarity shared-control
exemption). This is a one-time, structure-dependent reorder --- consistent
with the principle that local mixing belongs to the mixing stage, whose
job the gadget stage only sets up.

\textbf{Sandwich column float} (\code{9323272a}). In the sliced sandwich,
each middle-column \CNOT{} is assigned a registered random direction and
floats that way to its commutation extreme, dissolving the $C|N|D$ seam
--- the sandwich's most structure-revealing feature --- into a band before
gadgetization.

\section{Open items}

The production $n{=}128$ A/B (old vs.\ revised recipe through the
two-phase mixing pipeline, evaluated with \code{hmap\_affine} at degrees 1
and 2) has not yet been run. The persistent degree-2 snapshot legibility
of the linear XOR sharing remains the deepest open design question; the
candidate answers are a nonlinear share encoding (snapshots stop being
degree-1 in the wires) and generalized nonlinear RG3 injections
(compensated pair-writes with arbitrary conjunction fragments).

\appendix
\section{Gate-level diagrams of the CG menu}\label{app:cgfigs}

Each realization of Section~4 is drawn below in standard reversible-gate
notation. The five horizontal lines are the carrier wires (top to bottom:
$a^\ast$ the targeted carrier of $A$, then $b_0, b_1, c_0, c_1$); the
untouched second carrier of $A$ is a mask and is omitted. A gate is a
vertical line joining its pins. A filled pin~$\bullet$ is a positive
control (active when its wire is $1$), an open pin~$\circ$ a negative
control (active when $0$), and $\oplus$ the target that is flipped when the
gate is active. A plain conjunction gate is active exactly when all its
controls are satisfied. A \textbf{g57} gate carries a bar over its
target,~$\overline{\oplus}$: it is \emph{complemented}, active on the
complement of its control cube --- realizing $t \pe (x \vee \overline{y})$
from pins $\circ$ on $x$ and $\bullet$ on $y$. Columns are numbered to match
the gate sequences in Section~4 and apply left to right.

\bigskip
\begin{center}
\begin{minipage}[t]{0.47\textwidth}\centering
\textbf{0 --- collapse both}\\[4pt]
\begin{cgcirc}
\wires{5}
\gcol{1}{-1}{-2}\tg{1}{-1}\pc{1}{-2}
\gcol{2}{-3}{-4}\tg{2}{-3}\pc{2}{-4}
\gcol{3}{0}{-3}\tgb{3}{0}\nc{3}{-1}\pc{3}{-3}
\gcol{4}{-3}{-4}\tg{4}{-3}\pc{4}{-4}
\gcol{5}{-1}{-2}\tg{5}{-1}\pc{5}{-2}
\end{cgcirc}
\end{minipage}\hfill
\begin{minipage}[t]{0.47\textwidth}\centering
\textbf{1 --- collapse $c$}\\[4pt]
\begin{cgcirc}
\wires{5}
\gcol{1}{-3}{-4}\tg{1}{-3}\pc{1}{-4}
\gcol{2}{0}{-3}\tgb{2}{0}\nc{2}{-1}\pc{2}{-3}
\gcol{3}{0}{-3}\tgb{3}{0}\nc{3}{-2}\pc{3}{-3}
\gcol{4}{0}{-3}\tg{4}{0}\nc{4}{-3}
\gcol{5}{-3}{-4}\tg{5}{-3}\pc{5}{-4}
\end{cgcirc}
\end{minipage}

\vspace{1.6em}
\begin{minipage}[t]{0.47\textwidth}\centering
\textbf{2 --- collapse $b$}\\[4pt]
\begin{cgcirc}
\wires{4}
\gcol{1}{-1}{-2}\tg{1}{-1}\pc{1}{-2}
\gcol{2}{0}{-3}\tgb{2}{0}\nc{2}{-1}\pc{2}{-3}
\gcol{3}{0}{-4}\tg{3}{0}\nc{3}{-1}\pc{3}{-4}
\gcol{4}{-1}{-2}\tg{4}{-1}\pc{4}{-2}
\end{cgcirc}
\end{minipage}\hfill
\begin{minipage}[t]{0.47\textwidth}\centering
\textbf{3 --- linear tail}\\[4pt]
\begin{cgcirc}
\wires{6}
\gcol{1}{0}{-3}\tg{1}{0}\nc{1}{-3}
\gcol{2}{0}{-4}\tg{2}{0}\pc{2}{-4}
\gcol{3}{0}{-3}\tgb{3}{0}\nc{3}{-1}\pc{3}{-3}
\gcol{4}{0}{-3}\tgb{4}{0}\nc{4}{-2}\pc{4}{-3}
\gcol{5}{0}{-4}\tgb{5}{0}\nc{5}{-1}\pc{5}{-4}
\gcol{6}{0}{-4}\tgb{6}{0}\nc{6}{-2}\pc{6}{-4}
\end{cgcirc}
\end{minipage}

\vspace{1.6em}
\begin{minipage}[t]{0.47\textwidth}\centering
\textbf{4 --- legacy GADGET}\\[4pt]
\begin{cgcirc}
\wires{6}
\gcol{1}{-1}{-4}\tgb{1}{-4}\nc{1}{-1}\pc{1}{-2}
\gcol{2}{0}{-4}\tgb{2}{0}\pc{2}{-2}\nc{2}{-4}
\gcol{3}{0}{-4}\tgb{3}{0}\nc{3}{-1}\pc{3}{-4}
\gcol{4}{-1}{-4}\tgb{4}{-4}\nc{4}{-1}\pc{4}{-2}
\gcol{5}{0}{-3}\tgb{5}{0}\nc{5}{-2}\pc{5}{-3}
\gcol{6}{0}{-3}\tgb{6}{0}\pc{6}{-1}\nc{6}{-3}
\end{cgcirc}
\end{minipage}\hfill
\begin{minipage}[t]{0.47\textwidth}\centering
\textbf{5 --- quad $+$ pair}\\[4pt]
\begin{cgcirc}
\wires{6}
\gcol{1}{0}{-3}\tgb{1}{0}\nc{1}{-1}\pc{1}{-3}
\gcol{2}{0}{-3}\tgb{2}{0}\nc{2}{-2}\pc{2}{-3}
\gcol{3}{0}{-4}\tgb{3}{0}\nc{3}{-1}\pc{3}{-4}
\gcol{4}{0}{-4}\tgb{4}{0}\nc{4}{-2}\pc{4}{-4}
\gcol{5}{0}{-4}\tg{5}{0}\nc{5}{-3}\pc{5}{-4}
\gcol{6}{0}{-4}\tgb{6}{0}\nc{6}{-4}\pc{6}{-3}
\end{cgcirc}
\end{minipage}

\vspace{1.6em}
\begin{minipage}[t]{0.47\textwidth}\centering
\textbf{6 --- 4-cube ESOP}\\[4pt]
\begin{cgcirc}
\wires{4}
\gcol{1}{0}{-4}\tg{1}{0}\pc{1}{-1}\nc{1}{-4}
\gcol{2}{0}{-4}\tg{2}{0}\nc{2}{-2}\pc{2}{-4}
\gcol{3}{0}{-3}\tg{3}{0}\pc{3}{-2}\pc{3}{-3}
\gcol{4}{0}{-3}\tg{4}{0}\nc{4}{-1}\nc{4}{-3}
\end{cgcirc}
\end{minipage}
\end{center}

\bigskip\noindent
Reading these against the algebra: in 0--2 a \CNOT{} pair collapses a value
onto one carrier (the $\bullet$ on $b_1$/$c_1$ writing $b_0$/$c_0$) and the
mirror pair at the end restores it; variant~4 uses $c_1$ as scratch,
written and restored by the identical gates in columns~1 and~4. The
masking-safe realizations (3 and~6) are exactly those whose every gate
reads at most one carrier of each logical value.

\end{document}
